%DS12_Trang
%Đề cũ số 12. Đề thường: chỉ đáp án; câu khó có hướng dẫn ngắn.

\section{ĐỀ SỐ \thesection }
\setcounter{bt}{0}

%% Dạng 1. Thực hiện phép tính — 4 ý: a cơ bản, b tính hợp lý, c và d nhiều ngoặc và lũy thừa.
\begin{bt} Thực hiện phép tính:
	\begin{listEX}
		\item $33+(-47)+9$;
		\item $(-12)\cdot 25+(-12)\cdot 15+(-12)\cdot 10$;
		\item $\left\{2^{5}-\left[18-(-3)^{2}\right]\right\}\cdot(-2)+4^{0}$;
		\item $200-8\cdot\left[5^{2}-3\cdot\left(2^{3}-6^{0}\right)\right]$.
	\end{listEX}
	\loigiai{a) $-5$ \hfill b) $-600$ \hfill c) $-45$ \hfill d) $168$.}
\end{bt}

%% Dạng 2. Tìm x — 4 ý: a cơ bản, b trị tuyệt đối, c nhiều ngoặc, d có x ở cả hai vế.
\begin{bt} Tìm số nguyên $x$, biết:
	\begin{listEX}
		\item $4x=-28$;
		\item $|3x-9|=6$;
		\item $2(3x-4)-5(x-2)=7$;
		\item $3(2x-1)-x=4(x-2)+5$.
	\end{listEX}
	\loigiai{a) $x=-7$ \hfill b) $x=5$ hoặc $x=1$ \hfill c) $x=5$ \hfill d) $x=0$.}
\end{bt}

%% Dạng 3. Tập hợp — 2 ý, số nguyên và số tự nhiên.
\begin{bt} Biết số $0$ là số tự nhiên. Gọi $H$ là tập hợp các số nguyên $x$ thỏa mãn $-8\le x<7$ và $|x|$ chia hết cho $3$. Gọi $K$ là tập hợp các số tự nhiên nhỏ hơn $12$, không chia hết cho $2$ và chia hết cho $3$.
	\begin{enumerate}
		\item Liệt kê các phần tử của $H$.
		\item Liệt kê các số vừa thuộc $H$ vừa thuộc $K$. Liệt kê các số thuộc $H$ nhưng không phải số tự nhiên.
	\end{enumerate}
	\loigiai{a) $H=\{-6;\,-3;\,0;\,3;\,6\}$ \hfill b) Phần chung: $\{3\}$; thuộc $H$ nhưng không phải số tự nhiên: $\{-6;\,-3\}$.}
\end{bt}

%% Dạng 4. Ước và bội — thực tế chia hoa, học sinh tự nhận ra ƯCLN.
\begin{bt} Cô Lan có $42$ bông hồng và $28$ bông cúc. Cô cắm hết số hoa vào các lọ sao cho các lọ có số hồng bằng nhau, số cúc bằng nhau, và mỗi lọ có cả hai loại hoa.
	\begin{enumerate}
		\item Nếu số lọ lớn hơn $1$, cách cắm nào khiến mỗi lọ có nhiều hoa nhất? Khi đó có bao nhiêu lọ, mỗi lọ có bao nhiêu bông hồng và bao nhiêu bông cúc?
		\item Nếu muốn có nhiều lọ nhất thì có bao nhiêu lọ? Mỗi lọ có bao nhiêu bông hồng và bao nhiêu bông cúc?
	\end{enumerate}
	\loigiai{a) $2$ lọ; mỗi lọ $21$ bông hồng và $14$ bông cúc \hfill b) $14$ lọ; mỗi lọ $3$ bông hồng và $2$ bông cúc.}
\end{bt}

%% Dạng 5. Toán thực tế — chia đều số lượng, một đáp án trong khoảng cho trước.
\begin{bt} Một kho phát khẩu trang. Chia đều cho $20$ người hoặc cho $30$ người thì đều hết số khẩu trang, không dư chiếc nào. Biết trong kho có nhiều hơn $250$ chiếc và ít hơn $320$ chiếc. Hỏi kho có bao nhiêu chiếc khẩu trang?
	\loigiai{$300$ chiếc.}
\end{bt}

%% Dạng 6. Hình học thực tế — mảnh đất hình chữ L.
\begin{bt}
	\immini[thm]{
		Từ một mảnh đất hình chữ nhật dài $20$ m, rộng $12$ m, người ta cắt đi một hình chữ nhật $8$ m $\times$ $5$ m ở một góc để làm lối ra. Phần còn lại dùng để trồng rau, như hình vẽ.
		\begin{enumerate}
			\item Tính diện tích đất trồng rau.
			\item Tính chu vi phần đất trồng rau.
			\item Cứ mỗi mét vuông thu hoạch được $5$ kg rau. Giá bán là $9\,000$ đồng/kg. Tính số tiền thu được khi bán hết số rau.
		\end{enumerate}
	}{
		\begin{tikzpicture}[scale=0.28,line width=0.8pt,line join=round,>=stealth]
			\draw (0,0)--(12,0)--(12,5)--(20,5)--(20,12)--(0,12)--cycle;
			\node at (8,8) {Trồng rau};
			\draw[<->] (0,-1.2)--(12,-1.2) node[midway,below]{$12$ m};
			\draw[<->] (12.7,0)--(12.7,5) node[midway,right]{$5$ m};
			\draw[<->] (12,4.2)--(20,4.2) node[midway,below]{$8$ m};
			\draw[<->] (20.8,5)--(20.8,12) node[midway,right]{$7$ m};
			\draw[<->] (0,12.8)--(20,12.8) node[midway,above]{$20$ m};
			\draw[<->] (-1.4,0)--(-1.4,12) node[midway,left]{$12$ m};
		\end{tikzpicture}
	}
	\loigiai{a) $200$ m$^{2}$ \hfill b) $64$ m \hfill c) $9\,000\,000$ đồng.}
\end{bt}

%% Dạng 7. Nâng cao — so sánh lũy thừa bằng cách đưa về cùng cơ số.
\begin{bt} So sánh $2^{20}$ và $8^{6}$.
	\loigiai{Hướng dẫn: viết $8^{6}=(2^{3})^{6}=2^{18}$.\\
		Vì $20>18$ nên $2^{20}>2^{18}$. Vậy $2^{20}>8^{6}$.}
\end{bt}
